\scanpage{40}
$\mathcal X$ ens ordonné fini. On suppose que \deleted{\unsure{$\mathcal X$|il}}\inserted{$\bar{\mathcal X}=\mathcal X\cup$ plus petit él.} satisfait la condition

$(\Sigma)$ Soient $x\in\bar{\mathcal X}$, $y$ un successeur de $x$, $z$ un successeur de $y$. Alors il existe un unique $y'\in \mathcal X$, tel que $y'\ne y$, $x<y'<z$ ; \deleted{(de plus \unsure{on a} $y'$ est un successeur de $x$ et $z$ un successeur de $y'$}\deleted{\unsure{De m\textsuperscript{me}, si} $y$ est un él. minimal de $\mathcal X$, $z$ un successeur, $\exists! y'\in \mathcal X$ tel que $y'<z$ et $y'$ est minimal et}

Le $y'$ en question (noté $\sigma_{x,z}(y)$) est un successeur de $x$, et $z$ en est un successeur.

\diagram{\begin{tikzcd}[ampersand replacement=\&,row sep=tiny,column sep=tiny]
\&\deleted{\bullet}\arrow[dr]\&\\
\&\&\deleted{\bullet}\arrow[dr]\\
\deleted{\bullet}\arrow[uur]\arrow[r]\&\deleted{\bullet}\arrow[rr]\&\&\deleted{\bullet}
\end{tikzcd}}

En effet, si on avait \inserted{p.ex.} $x<y''<y'$, on aurait $x<y''<z$, donc (comme $y''\ne y'$) $y''=y$ donc $y<y'<z$, contrairement à l'hyp. que $z$ est un succ. de $y$.

\diagram{\begin{tikzcd}[ampersand replacement=\&,row sep=tiny,column sep=small]
\deleted{x<y<}\&\&\\[-2pt]
\&y'\arrow[dr]\&\\
x\arrow[ur]\arrow[dr]\&\&z\\
\&y\arrow[ur]\&
\end{tikzcd}}

% margin note written sideways
\margin{\uline{On suppose} Pour $x\in \mathcal X$ les chaines \unsure{ascendantes} \inserted{vers $x$} maximales \unsure{ont} \unsure{ttes} \unsure{$=$|m\textsuperscript{me}} longueur.}
$n$ répère de $\mathcal X$ \inserted{\unsure{($n\ne0$)|($n\ge0$)}} : suite
\[
x_0,x_1,\ldots,x_n
\]
d'éléments de $\mathcal X$, tels que a) $x_0$ est minimal de $\mathcal X$

b) $\forall\ 0\le i\le n-1$, $x_{i+1}$ est un successeur de $x_i$. Soit $R_n(\mathcal X)$ l'ens des $n$-répères

\uline{Structures}

a) $R_n\xrightarrow{p_n}R_{n-1}$ \qquad $(n\ge1)$

b) Opérations $\sigma_0,\sigma_1,\ldots,\sigma_{n-1}$ \quad (si $n\ge1$)

\uline{Propriétés} \quad a) les $\sigma_i$ sont des involutions $(n\ge1)$

b) Soit $n\ge3$, et $0\le i<j\le n-1$, avec $j\ne i+1$, alors $\sigma_i$ et $\sigma_j$ commutent \deleted{\unsure{et}} (dans $R_n$).

c) Soit $0\le i_1<i_2<\ldots<i_\nu\le n-1$ \inserted{($\nu\ge1$)}, avec $i_{\alpha+1}\ne i_\alpha+1$ $\forall\alpha=1,\ldots,\nu-1$. Alors $\sigma_{i_1}\sigma_{i_2}\ldots\sigma_{i_\nu}$ (qui est une involution de $R_n$) est sans pt fixe.

d) $R_n\xrightarrow{p_n\cdots p_{p+1}}R_p$ commute à $\sigma_{n-1}\ldots\sigma_{p+1}$ et $\sigma_{p-1},\ldots\sigma_0$

% archive sheet number: 39
